\medskip

		\begin{enumerate}
			\item Soit $f(x)=k$ une fonction constante définie sur $\R$ solution de $\left( E_0\right)$.

On a donc $f'=f$ soit $0=k$.

L'unique fonction constante solution de l'équation différentielle $\left(E_{0}\right)$ est donc la fonction nulle.

			\item Les solutions de l'équation différentielle $\left(E_{0}\right)$ sont les fonctions de la forme 

$f(x)=C\e^x$ avec $C \in \R$.



		\item


Pour tout réel $x$ on a : \quad $h'(x) =-2\sin(x) +\cos(x)$.

D'autre part : \quad $h(x)-\cos (x)-3 \sin (x)= 2 \cos (x)+\sin (x)- \cos (x)-3 \sin (x) =$

$ \cos(x)-2\sin(x)$

d'où pour tout réel $x$,  \quad$h'(x) = h(x) -\cos (x)-3 \sin (x)$,  c'est à dire, $h$ est solution de l'équation différentielle $(E)$.

			\item Supposons que  $f$ soit une solution de $(E)$.

Pour tout réel $x$ on a :

$\begin{aligned}[t] (f-h)'(x)=f'(x)-h'(x) &\implies (f-h)'(x)= f(x) -\cos (x)-3 \sin (x) - \left( h(x) -\cos (x)-3 \sin (x) \right)\\
&\qquad \text{  car $f$ et $h$ sont solutions de $(E)$}\\
&  \Longrightarrow (f-h)'(x)= f(x) -\cos (x)-3 \sin (x) - h(x) +\cos (x)+3 \sin (x) \\
&  \Longrightarrow (f-h)'(x)= f(x)  - h(x)   = (f-h)(x)\end{aligned}$

Donc $f-h$ est solution de $\left(E_{0}\right)$
\smallskip

Réciproquement : supposons que $f-h$ soit solution de $\left(E_{0}\right)$

On a donc $(f-h)'(x) = f(x)  - h(x)  $ soit $f'(x) -h'(x)= f(x) -h(x)$

D'où : $f'(x) = f(x) -h(x) +h'(x) = f(x) -h(x) + h(x) -\cos (x)-3 \sin (x)$ car $h$ est solution de $(E)$

Donc : $f'(x) =  f(x) -\cos (x)-3 \sin (x)$ c'est à dire $f$ est solution de $(E)$.

\emph{Conclusion :} $f$ est solution de $(E)$ si et seulement si $f-h$ est solution de $(E_0)$.
		\item

On a donc $f(x)-h(x)= C\e^x$ avec $C \in \R$

Toutes les solutions de l'équation différentielle $(E)$ sont donc les fonctions

$f(x)=C\e^x +2\cos(x)+\sin(x)$ avec $C \in \R$

\item $g$ est solution de l'équation différentielle $(E)$  donc il existe un réel $C$ tel que

$g(x)=C\e^x +2\cos(x)+\sin(x)$.

De plus  $g(0)=0$ donc $g(0)=C\e^0 +2\cos(0)+\sin(0)=0$.

D'où $C\times 1 +2\times 1 + 0 =0 \Longleftrightarrow C=-2$

On a donc :\quad $g(x)=-2\e^x +2\cos(x)+\sin(x)$.

		\item
Calculons : \quad$\displaystyle
I = \int_{0}^{\frac{\pi}{2}}\left(-2 \e^{x}+\sin (x)+2 \cos (x)\right)\,  \mathrm{d} x = \int_{0}^{\frac{\pi}{2}} g(x)\,  \mathrm{d} x $


$\displaystyle
I = 	\int_{0}^{\frac{\pi}{2}}\left(-2 \e^{x}+\sin (x)+2 \cos (x)\right)\,  \mathrm{d} x = \Big[ -2\e^{x} -\cos(x) +2\sin(x) \Big]_0^ {\frac{\pi}{2}}$

$I= -2\e^{\frac{\pi}{2}} -\cos\left(  {\dfrac{\pi}{2}}\right) +2\sin\left(  {\dfrac{\pi}{2}}\right)  - \left(  -2\e^{0} -\cos(0) +2\sin(0)  \right) $

$I= -2\e^{\frac{\pi}{2}} -0+2\times 1 - \left( -2 -1+2\times 0\right) =  -2\e^{\frac{\pi}{2}}+ 2 +3  = -2\e^{\frac{\pi}{2}}+ 5  $

		\end{enumerate}

\medskip
